% DSAIT4335 Recommender Systems -- Final Project report (hand-in version)
%
% Formatting required by the brief: Times New Roman, 12 pt, line spacing 1.15.
% Structure required by the brief: first page = group information; main body = at most ONE page per task
% (results + discussion) with at most 200 WORDS of discussion per task; appendix = additional results
% (same limits). The main body must be self-contained.
%
% Who writes what (RecSys-Project-Plan-and-Team-Split, Part 2): each member writes the subsections of the
% sub-tasks they own; A edits the whole report for one voice and the page limits.
%   A  Models        1.1, 1.2, 2.2        C  Weighted hybrid  1.3, 1.5 (weighted), 2.3, 2.6
%   B  Evaluation    2.1, 2.5, 3.4        D  Content-based    1.1 (content-based), 1.4, 1.5 (others), 2.4, 3.3
%   E  Re-rankers    3.1, 3.2
%
% Build: cd report && tectonic report.tex      (or: latexmk -pdf -outdir=build report.tex)
% Overleaf: Menu -> Main document -> report/report.tex
%
% Generated assets (committed to git, produced by python -m project.experiments.<...>):
%   ../figures/generated/<name>.pdf        -> \generatedfigure[width=...]{<name>}
%   tables/generated/<name>.tex            -> \generatedtable{<name>}
% A red box is shown for an asset that has not been generated/committed yet, so the report always compiles.
\documentclass[12pt,a4paper]{article}

\usepackage[T1]{fontenc}
\usepackage[utf8]{inputenc}
\usepackage{newtxtext,newtxmath}          % Times New Roman (text + maths)
\usepackage{setspace}
\setstretch{1.15}                         % line spacing 1.15
\usepackage[a4paper,margin=2.2cm]{geometry}
\usepackage{graphicx}
\usepackage{booktabs}
\usepackage{amsmath,amssymb}
\usepackage{xcolor}
\usepackage[font=small,skip=4pt]{caption}
\usepackage{enumitem}
\setlist{nosep,leftmargin=*}
\usepackage[hidelinks]{hyperref}
\usepackage{titlesec}
\titlespacing*{\section}{0pt}{8pt}{4pt}
\titlespacing*{\subsection}{0pt}{6pt}{2pt}

% ---------------------------------------------------------------------------
% Title-page data -- fill in
% ---------------------------------------------------------------------------
\newcommand{\coursecode}{DSAIT4335}
\newcommand{\coursename}{Recommender Systems}
\newcommand{\projecttitle}{Hybrid Recommendation on MovieLens-100K:\\ Accuracy, Coefficients and Societal Re-ranking}
\newcommand{\groupnumber}{\TODO{group number}}

% ---------------------------------------------------------------------------
% Helpers
% ---------------------------------------------------------------------------
\newcommand{\TODO}[1]{\textcolor{red}{\textbf{[TODO: #1]}}}
\newcommand{\owner}[1]{\marginpar{\footnotesize\color{gray}#1}}   % who writes this subsection (remove before hand-in)

\newcommand{\missingasset}[1]{%
  \fbox{\parbox{0.85\linewidth}{\centering\color{red}%
    \textbf{TODO:} generated asset not found\\ \texttt{#1}}}}

\graphicspath{{../figures/generated/}{figures/}}
\newcommand{\generatedfigure}[2][width=\linewidth]{%
  \IfFileExists{../figures/generated/#2.pdf}%
    {\includegraphics[#1]{../figures/generated/#2.pdf}}%
    {\IfFileExists{figures/generated/#2.pdf}%
      {\includegraphics[#1]{figures/generated/#2.pdf}}%
      {\missingasset{figures/generated/\detokenize{#2}.pdf}}}}
\newcommand{\generatedtable}[1]{%
  \IfFileExists{tables/generated/#1.tex}%
    {\input{tables/generated/#1.tex}}%
    {\IfFileExists{report/tables/generated/#1.tex}%
      {\input{report/tables/generated/#1.tex}}%
      {\missingasset{report/tables/generated/\detokenize{#1}.tex}}}}

\begin{document}

% ===========================================================================
% First page: group information
% ===========================================================================
\pagenumbering{roman}
\begin{titlepage}
  \centering
  \vspace*{3cm}
  {\large \coursecode\ \coursename\ --- Final Project\par}
  \vspace{1.5cm}
  {\LARGE\bfseries \projecttitle\par}
  \vspace{1.5cm}
  {\Large Group \groupnumber\par}
  \vspace{1.5cm}
  {\large\bfseries Group members\par}
  \vspace{0.5em}
  \begin{tabular}{ll}
    \TODO{name} & \TODO{student number}\\
    \TODO{name} & \TODO{student number}\\
    \TODO{name} & \TODO{student number}\\
    \TODO{name} & \TODO{student number}\\
    \TODO{name} & \TODO{student number}\\
  \end{tabular}
  \vfill
  {\large Delft University of Technology\par}
  {\normalsize Code: \texttt{project/} in the submitted archive (see \texttt{project/models/README.md} for how to run it)\par}
\end{titlepage}
\pagenumbering{arabic}

% ===========================================================================
% TASK 1 -- Hybrid recommender (max. one page, max. 200 words of discussion)
% ===========================================================================
\section{Task 1 --- Hybrid Recommender}
\label{sec:task1}

\subsection{Individual recommender models (Tasks 1.1 and 1.2)}\owner{A}
\label{sec:task1-individual}
\paragraph{Setup (shared by all tasks).}
All experiments use RecBole on MovieLens-100K with \emph{one} frozen split: seed 2020, a random
80/10/10 split per user (RecBole \texttt{RS}, \texttt{RO}, \texttt{group\_by: user}), full ranking over all
items the user has not interacted with, and every held-out interaction counted as relevant (RecBole's default; no
rating threshold). Train/validation/test hold 80\,808 / 9\,596 / 9\,596 interactions of 943 users and 1\,682
items; the split is exported once with a checksum manifest and verified to be identical through every model's
configuration. Every model exports its score for all user--item pairs and a top-50 candidate list per user
(validation lists mask the training history, test lists mask train\,+\,validation); the final list is the top 10.
Hybrids and re-rankers only read these files, so nothing is retrained downstream.

\paragraph{Individual models (Task 1.1).}
We use the course's RecBole models, one per family discussed in class, plus the two na\"ive baselines: Random
(scores drawn per user; RecBole's own \texttt{Random} gives all users of a batch the same list) and Pop;
neighbourhood models UserKNN and ItemKNN (cosine, RecBole \texttt{ItemKNN} with \texttt{knn\_method});
the linear item--item models EASE and SLIMElastic; BPR matrix factorisation (pairwise ranking loss);
the neural NeuMF and FISM; and the graph models LightGCN and NGCF. The content-based model of
Assignment~1 is added by Track~D on the same split.

\paragraph{Tuning (Task 1.2).}
Each model is tuned by an exhaustive grid (7--20 configurations, Table~\ref{tab:tuning-summary} in the appendix)
on \emph{validation} NDCG@10, the target agreed for every later choice in the project; the test set is used
once, for the selected configuration. Neural models train for up to 100 epochs with validation after every
epoch and early stopping after 10 epochs without improvement, instead of the course configs' fixed 20 epochs.
Selected: ItemKNN $k{=}200$, UserKNN $k{=}50$ (both without shrinkage), EASE $\lambda{=}500$, SLIM
$\alpha{=}0.1$, $l_1{=}0.01$, BPR 128 factors with learning rate $10^{-3}$, NeuMF MLP $[128,64,32]$ with
dropout 0.1, LightGCN 3 layers with learning rate $5\cdot10^{-3}$.

\emph{Discussion.} Tuning matters most where the course configuration left a model under-trained: LightGCN
rises from 0.136 to 0.243 validation NDCG@10 (its best epoch is 94 of 100) and BPR from 0.199 to 0.241, both
mainly through a higher learning rate or a longer schedule; the memory-based and linear models move by at most
0.003 because, with 85 interactions per user on average, item--item statistics are already well estimated
and extra neighbours or shrinkage barely change the ranking. EASE prefers strong regularisation
($\lambda{=}500$; $\lambda{=}10$ loses 0.05), consistent with a dense, small catalogue in which an
unregularised item--item solution overfits co-occurrence noise.

\TODO{D: one sentence on the content-based model (Assignment 1 BERT model on the same split, exported as a score file).}

\subsection{Weighted hybrid with learned coefficients (Task 1.3)}\owner{C}
\TODO{C: regression model used to learn the coefficients; fitted on the validation candidates (union of each model's top-50); per-user score normalisation; the learned weights.}

\subsection{Other hybrid models (Tasks 1.4 and 1.5)}\owner{D, C}
\TODO{D: switching and mixed hybrids and why; C/D: how all hybrids were tuned on validation NDCG@10.}

\clearpage
% ===========================================================================
% TASK 2 -- Evaluation of effectiveness (max. one page, max. 200 words of discussion)
% ===========================================================================
\section{Task 2 --- Evaluation of Effectiveness}
\label{sec:task2}

\subsection{Evaluation metrics (Task 2.1)}\owner{B}
\label{sec:task2-metrics}
All metrics are our own implementation (\texttt{project/metrics}); RecBole only produces the lists, so models,
hybrids and re-ranked lists are scored by identical code. Each list is cut at $K{=}10$, scored per user and
macro-averaged over the 943 users; catalogue-level metrics are computed once over all lists. Item popularity
$c_i$, user profiles and all groups come only from the training split (identical for validation and test), never
from held-out data. \emph{Accuracy:} Precision, Recall, NDCG, MRR and MAP. \emph{Beyond accuracy}
(Table~\ref{tab:beyond-accuracy}; formulas and our choices in Table~\ref{tab:metric-definitions}): intra-list
diversity (ILD, mean pairwise Jaccard distance between genre sets), catalogue coverage, novelty (mean
self-information $-\log_2$ of a film's share of interactions), serendipity (relevant films that a most-popular
list would not have shown), genre miscalibration MC$_{KL}$ between history and list, average popularity and
long-tail share; for fairness, the user-side GRU (largest NDCG gap between the low-, medium- and high-activity
users of Task~2.5) and popularity deviation UPD, and the item-side Gini index of exposure. Validation: on all 22 lists our accuracy
metrics equal the values RecBole reported to its 4-decimal rounding, and coverage, Gini, entropy and AvgPop equal
RecBole's own metric code applied to our lists (Table~\ref{tab:metric-validation}); the metrics RecBole lacks are
pinned down by hand-computed tests, including the lecture's worked examples.

\begin{table}[htbp]
  \centering
  \caption{Beyond-accuracy metrics on the test split (lists of 10; $\uparrow$/$\downarrow$: higher/lower is
  better). Nov.: novelty in bits; Ser.: serendipity; AvgPop: share of users who rated a recommended film; Tail:
  share of long-tail films; GRU: largest NDCG@10 gap between user-activity groups; Cov.\ and Gini: catalogue
  coverage and exposure inequality over all 1\,682 films.}
  \label{tab:beyond-accuracy}
  \small\setlength{\tabcolsep}{3pt}
  \generatedtable{beyond_accuracy_results}
\end{table}

\emph{Discussion.} The beyond-accuracy metrics separate the models by family more than accuracy does. The
tuned embedding models (BPR, NeuMF, LightGCN) cover 32--37\% of the catalogue, recommend less popular films (AvgPop
0.23) and follow each user's taste for popular or niche films best (UPD 0.14). The neighbourhood and
item--item linear models cover only 17--23\% (AvgPop 0.24--0.27): they score a film by aggregating its
co-occurrence with the user's films or neighbours, which favours films that co-occur with many others. EASE and
SLIM are the best calibrated (MC$_{KL}$ 0.69 and 0.66; all others $\geq0.74$), as they reconstruct each user's
history from item--item weights, so the recommended films resemble the user's own. Two properties of the metrics
matter for Task~3: genre ILD barely varies (Random's 0.82 equals the mean distance of two random films, since
65\% of film pairs share no genre), and GRU largely reflects the protocol: high-activity users hold 22 held-out films
on average against 2.5 for low-activity users, so even Random gives them four times the NDCG
(Section~\ref{sec:task2-groups}).

\subsection{Comparison with the baselines (Task 2.2)}\owner{A}
\label{sec:task2-baselines}
Table~\ref{tab:model-results} reports the tuned models on the test split (accuracy metrics computed by
RecBole on the exported lists; beyond-accuracy columns are added from Section~\ref{sec:task2} once Track~B's
metrics are in; FISM and NGCF are reported with the course configuration, untuned).

\begin{table}[htbp]
  \centering
  \caption{Test-set accuracy of the individual models (lists of 10, all held-out interactions relevant); best
  value per column in bold. The last column is the validation NDCG@10 used for model selection.}
  \label{tab:model-results}
  \small\setlength{\tabcolsep}{4pt}
  \generatedtable{model_results}
\end{table}

\emph{Discussion.} Pop is a strong na\"ive baseline here: one unpersonalised list hits a relevant film for
47\% of the users (Random: 7\%) because MovieLens-100K's interactions concentrate on a few hundred films.
All personalised models beat it, yet 47--61\% of their top-10 slots are films Pop also lists, and their
recommendations were rated by 23--27\% of the training users on average (Pop 33\%, Random under 5\%).
The item--item linear models EASE and SLIM are the most accurate: with 943 users and 1\,682 items the
co-occurrence matrix is dense enough for a closed-form, strongly regularised solution, whereas the
embedding models need long training to reach a similar level on this little data. The accuracy winners
thus inherit much of Pop's popularity bias, which Tasks~2.5 and~3 examine.

\subsection{Hybrid coefficient analysis (Task 2.3)}\owner{C}
\TODO{C: contribution of each individual recommender to the hybrid (coefficients, drop-one ablation).}

\subsection{Further analyses (Task 2.4)}\owner{D}
\TODO{D: overlap between the models' lists; what an oracle best-model-per-user would score.}

\subsection{User and item groups (Task 2.5)}\owner{B}
\label{sec:task2-groups}
Groups are built from the training split only and frozen for all models
(Table~\ref{tab:group-definitions}). Users are split twice, into tertiles: by \emph{activity} (number of training
interactions: low $\leq34$, medium $\leq92$, high) and by \emph{taste}, the share of head films in their training
interactions (niche $\leq19\%$, mixed $\leq29\%$, mainstream), the fairness lecture's users of unpopular, both, or
popular items. Films are head, mid or tail (20/60/20\% of the training interactions); 27 films without a training
interaction are kept apart. Metrics are computed on each complete list and averaged per group with 95\% bootstrap
intervals; two models are compared on the same users with paired bootstraps, on validation and test; item groups use the micro recall of
their held-out pairs and their share of recommendation slots (Tables~\ref{tab:user-groups} and~\ref{tab:item-groups}).

\begin{figure}[htbp]
  \centering
  \generatedfigure[width=\linewidth]{user_group_analysis}
  \caption{Test-set accuracy per group: dots are group means, whiskers 95\% bootstrap intervals, the grey line the
  spread between groups. (d) Recall of held-out films by popularity group (support: 1\,863 head, 5\,696 mid and
  2\,016 tail pairs).}
  \label{fig:user-groups}
\end{figure}

\emph{Discussion.} Raw group gaps mostly reflect the protocol: high-activity users hold 22 held-out films against
2.5 for low-activity ones, so NDCG@10 is highest and (except for Random) Recall@10 lowest for high-activity users
(Figure~\ref{fig:user-groups}a,b); niche users (18 held-out films) and mainstream users (4) show the same reversal.
We therefore compare models within groups and keep only differences that hold on both validation and test.
EASE's lead over the embedding models is clear for high-activity users (paired NDCG $+0.03$ to $+0.05$ over
LightGCN and BPR on both splits) but vanishes for low-activity users on validation: with few interactions no model
is reliably better, so switching models by activity has little to exploit (Task~2.6). On the item side, the 58
head films take 47--62\% of the tuned models' slots, while tail films, 21\% of the held-out pairs, are almost
never retrieved (recall $\leq0.021$). The embedding models expose and retrieve about ten times more tail films
than EASE, whereas Random fills 70\% of its slots with tail films and still retrieves fewer: exposure alone is not
success. Niche users are served worst in their taste: EASE and UserKNN give them 50--57\% head films against
15\% in their profiles (UPD 0.22--0.26 vs.\ 0.14--0.15 for mainstream users;
Figure~\ref{fig:popularity-alignment}), while BPR and LightGCN keep UPD at 0.13--0.15 in every taste group.

\TODO{B: add the hybrids once C and D export their lists (re-run \texttt{group\_analysis} with \texttt{--figure-names}).}

\subsection{Insights (Task 2.6)}\owner{C}
\TODO{C: group-specific hybrid motivated by 2.5.}

\clearpage
% ===========================================================================
% TASK 3 -- Societal aspects (max. one page, max. 200 words of discussion)
% ===========================================================================
\section{Task 3 --- Societal Aspects}
\label{sec:task3}

\subsection{Re-rankers (Task 3.1)}\owner{E}
\TODO{E: diversity (MMR), calibration, user-side and item-side fairness re-rankers on the top-50 candidates.}

\subsection{Effect of the re-rankers (Task 3.2)}\owner{E}
\TODO{E: trade-off curves (NDCG@10 against the target metric over $\lambda$) on the best models and the hybrid.}

\begin{table}[htbp]
  \centering
  \caption{Results with re-rankers applied.}
  \label{tab:reranker-results}
  \generatedtable{reranker_results}
\end{table}

\subsection{Hybrids combined with re-rankers (Task 3.3)}\owner{D}
\TODO{D: re-rank then combine versus combine then re-rank.}

\subsection{Impact on user and item groups (Task 3.4)}\owner{B}
\TODO{B: before/after re-ranking per group.}

\clearpage
% ===========================================================================
% APPENDIX -- additional results (same limits per task)
% ===========================================================================
\appendix
\section{Additional results for Task 1}
\label{app:task1}
\subsection{Hyper-parameter search (Task 1.2)}
Search spaces (\texttt{project/configs/hyper/}): ItemKNN and UserKNN $k\in\{10,20,50,100,200\}$,
shrink $\in\{0,10,50,100\}$; EASE $\lambda\in\{10,50,100,250,500,1000,2000\}$; SLIMElastic
$\alpha\in\{0.1,0.2,0.5,1\}$, $l_1\in\{0.01,0.02,0.05,0.1\}$; BPR factors $\in\{32,64,128\}$,
learning rate $\in\{5\cdot10^{-4},10^{-3},5\cdot10^{-3}\}$; NeuMF learning rate $\in\{10^{-3},5\cdot10^{-3}\}$,
MLP $\in\{[64,32],[128,64,32]\}$, dropout $\in\{0.1,0.3\}$; LightGCN layers $\in\{1,2,3\}$, learning rate
$\in\{10^{-3},5\cdot10^{-3}\}$, $L_2\in\{10^{-4},10^{-3}\}$. Every configuration is listed in
\texttt{results/processed/tuning/<Model>.csv}. Two selected values sit on the border of their grid
(ItemKNN $k{=}200$, BPR 128 factors); the neighbouring value is within 0.003 NDCG@10, so the grids were not
extended. All numbers are single runs with seed 2020; a 3-seed variance check is left for the final release.

\begin{table}[htbp]
  \centering
  \caption{Grid search per model: number of configurations, selected values, and NDCG@10 before (course
  configuration) and after tuning on validation and test.}
  \label{tab:tuning-summary}
  \small\setlength{\tabcolsep}{4pt}
  \generatedtable{tuning_summary}
\end{table}

\subsection{Accuracy overview (Task 2.2)}
\begin{figure}[htbp]
  \centering
  \generatedfigure[width=\linewidth]{model_comparison}
  \caption{Test-set NDCG@10, Recall@10 and Precision@10 per model (hatched: baselines; dotted line: Pop).}
  \label{fig:model-comparison}
\end{figure}

\section{Additional results for Task 2}
\label{app:task2}
\subsection{Metric definitions, conventions and validation (Task 2.1)}
\label{app:metrics}
Table~\ref{tab:metric-definitions} lists every metric of \texttt{project/metrics} (formulas, edge cases and usage
in \texttt{project/metrics/README.md}). Per-user values are macro-averaged over the evaluated users (every user
with at least one held-out item: all 943); coverage and Gini are computed once over all lists, never per batch.
\emph{Groups} (the groups of Task~2.5, Table~\ref{tab:group-definitions}) are built from the training split:
films sorted by $c_i$, the head holds the most popular films that together account for 20\% of the training
interactions (58 films, $c_i\geq216$), the tail the least popular films that together account for 20\% (1\,118
films, $1\leq c_i\leq48$), the mid the remaining 479; the 27 films without a training interaction form a
separate unseen group (long tail for TailShare). \emph{Edge cases:} a list with a duplicate, an unknown film or a film from the user's history is
rejected; a missing list counts as empty (accuracy 0); undefined values (ILD of fewer than two films) are left
out of the mean and counted. None of these occurs for the exported lists.

\begin{table}[htbp]
  \centering
  \caption{Our metrics ($K{=}10$; $L_u$: list, $T_u$: held-out relevant films, $H_u$: training interactions of user
  $u$, $c_i$: training interactions of film $i$, $I$: the 1\,682 films). ``Choice'' marks a convention the lectures leave open.}
  \label{tab:metric-definitions}
  \small\setlength{\tabcolsep}{4pt}
  \begin{tabular}{@{}p{2.3cm}p{9.1cm}p{4.3cm}@{}}
    \toprule
    Metric & Definition & Lecture; our choice \\
    \midrule
    Precision, Recall & $|L_u\cap T_u|/K$;\; $|L_u\cap T_u|/|T_u|$ & Evaluation 1 \\
    NDCG & $\sum_{r\leq K}\mathrm{rel}_r/\log_2(r{+}1)$, divided by its ideal value with $\min(K,|T_u|)$ hits & Evaluation 2 \\
    MRR & $1/$rank of the first relevant film, 0 if none & Evaluation 2 \\
    MAP & $\sum_{r\leq K}\mathrm{P@}r\cdot\mathrm{rel}_r/\min(K,|T_u|)$ & Evaluation 2; denominator$^{a}$ \\
    ILD & mean of $d(i,j)=1-|G_i\cap G_j|/|G_i\cup G_j|$ over all film pairs of $L_u$ ($G_i$: genres) & Diversity; Jaccard on genres \\
    Coverage & $|\bigcup_u L_u|/|I|$ (global) & Diversity; Fairness \\
    Novelty & mean over $L_u$ of $-\log_2\bigl(c_i/\sum_j c_j\bigr)$, unseen films counted as $c_i{=}1$ & Diversity \\
    Serendipity & $|\{i\in L_u\cap T_u:\ i\notin \mathrm{Pop}_u\}|/K$, $\mathrm{Pop}_u$: the $K$ most popular films the user has not interacted with & Diversity; formula \\
    MC$_{KL}$ & $\sum_g p(g|u)\ln\frac{p(g|u)}{(1-\alpha)\,q(g|u)+\alpha\, p(g|u)}$, $\alpha{=}0.01$; $p$ from $H_u$, $q$ from $L_u$, a film's $m$ genres weighted $1/m$ & Calibration; uniform weights \\
    AvgPop & mean over $L_u$ of $c_i/|U|$ (share of users who rated the film) & popularity bias \\
    TailShare & share of $L_u$ in the tail or unseen group & groups (above) \\
    UPD & $\mathrm{JSD}\bigl(P(H_u),P(L_u)\bigr)$, $P$: shares of head/mid/tail/unseen films, log base 2 & Fairness; $P$ and JSD \\
    GRU & $\max_g\overline{\mathrm{NDCG}}_g-\min_g\overline{\mathrm{NDCG}}_g$ over the low/medium/high activity groups (for two groups the lecture's $|\overline{\mathrm{NDCG}}_{1}-\overline{\mathrm{NDCG}}_{2}|$) & Fairness; groups, max$-$min \\
    Gini & $\sum_j(2j-n-1)\,e_{(j)}/(n\sum_j e_j)$, $e_i$: number of lists containing film $i$, all $n{=}|I|$ films (global) & Fairness \\
    \bottomrule
  \end{tabular}\\[2pt]
  {\footnotesize $^{a}$The lecture's AP divides by the ``number of relevant items'', but its worked example divides
  by the number of hits; we use $\min(K,|T_u|)$, as RecBole does.}
\end{table}

\begin{table}[htbp]
  \centering
  \caption{Validation on all 22 exported lists (11 models, validation and test split): largest absolute difference
  between our value and RecBole's, either the value RecBole reported for the same checkpoint (rounded to 4 decimals)
  or RecBole's metric code applied to our lists. Entropy is the normalised exposure entropy.}
  \label{tab:metric-validation}
  \small
  \generatedtable{metric_validation}
\end{table}

\subsection{User and item groups (Task 2.5)}
\label{app:groups}
Activity and taste are related (Pearson $r=-0.53$ between training interactions and head share): 62\% of the
niche users are highly active, so the taste and activity effects in Figure~\ref{fig:user-groups} are partly
the same users. The activity $\times$ taste cells, the per-group values of every metric with their intervals and the
paired differences between every two models are in \texttt{results/processed/user\_group\_*\_test.csv}; the
validation-split versions are the ones to use for design decisions (Task~2.6).

\begin{table}[htbp]
  \centering
  \caption{Groups of Task~2.5 (training split only; tertile cuts are value thresholds, so equal values share a group).}
  \label{tab:group-definitions}
  \small
  \generatedtable{group_definitions}
\end{table}

\begin{table}[htbp]
  \centering
  \caption{Test-set metrics per user group (means over the group's users). UPD: popularity deviation between a
  user's training interactions and list (lower is better).}
  \label{tab:user-groups}
  \small\setlength{\tabcolsep}{2.5pt}
  \generatedtable{user_group_results}
\end{table}

\begin{table}[htbp]
  \centering
  \caption{Test-set recall of the held-out films of each popularity group (micro: all held-out pairs of the group
  count once; support 1\,863 head, 5\,696 mid, 2\,016 tail pairs) and share of all recommendation slots. The 21
  held-out pairs on unseen films are retrieved by no model; only Random recommends unseen films (1.7\% of slots).}
  \label{tab:item-groups}
  \small\setlength{\tabcolsep}{4pt}
  \generatedtable{item_group_results}
\end{table}

\begin{figure}[htbp]
  \centering
  \generatedfigure[width=\linewidth]{popularity_alignment}
  \caption{Popularity mix of the users' training interactions (top bar) and of each model's test lists, per taste
  group: every personalised model recommends more head films than the profiles contain, relative to the profile
  most of all for niche users; the embedding models stay closest to the profiles.}
  \label{fig:popularity-alignment}
\end{figure}

\TODO{B, C, D: full per-group tables, extra metrics, alternative relevance rule.}

\section{Additional results for Task 3}
\TODO{E: extra $\lambda$ values, per-group tables after re-ranking.}

\end{document}
